\documentclass[11pt]{article}

\usepackage[T1]{fontenc}
\usepackage[margin=1in]{geometry}
\usepackage{amsmath,amssymb,amsthm,mathtools}
\usepackage{graphicx}
\usepackage{microtype}
\usepackage[hidelinks]{hyperref}

\newtheorem{theorem}{Theorem}[section]
\newtheorem{proposition}[theorem]{Proposition}
\newtheorem{lemma}[theorem]{Lemma}
\newtheorem{corollary}[theorem]{Corollary}
\theoremstyle{remark}
\newtheorem{remark}[theorem]{Remark}

\newcommand{\RNP}{\mathrm{RNP}}
\newcommand{\GRNP}{G\text{-}\mathrm{RNP}}
\newcommand{\Homeo}{\operatorname{Homeo}}
\newcommand{\Iso}{\operatorname{Iso}}
\newcommand{\supp}{\operatorname{supp}}
\newcommand{\co}{\operatorname{co}}

\title{A Counterexample to \texorpdfstring{$G$-RNP $\Rightarrow$ RNP}{G-RNP implies RNP}\\
\large Candidate counterexample packet for Problem 8.2}
\author{Automatic Math Discovery Run \texttt{fa\_banach\_001}}
\date{June 22, 2026}

\begin{document}
\maketitle

\begin{abstract}
Dantas, Doucha, Jung and Raunig asked whether the group-equivariant
Radon--Nikodym property implies the classical Radon--Nikodym property.
The supplied finite-group proof is correct.  Its induction idea can in fact
be made to work for every locally compact $\sigma$-compact group.  There is
also an elementary observation, apparently missed in the source paper, that
its printed notion of weak $G$-dentability is exactly ordinary dentability.
The unrestricted implication nevertheless fails.  Let
$G=\Homeo_+([0,1])$ and let it act by composition on
$X=C_0((0,1))$.  A theorem of Megrelishvili says that every continuous
isometric representation of $G$ on a reflexive Banach space is trivial.
The Mazur map then shows that every continuous measure-algebra action
allowed in the definition of the $G$-RNP is trivial.  Since $X^G=\{0\}$,
$X$ has the $G$-RNP, while $X$ fails the classical RNP because it contains
an isometric copy of $c_0$.  Hence $G$-RNP $\Rightarrow$ RNP is false for
general topological groups.
\end{abstract}

\section{Packet status and source evidence}

This packet upgrades and absorbs the earlier finite-group partial packet,
which is preserved in the local history directory.  The supplied TeX/PDF
report is stored in a dated evidence directory.

The source target is Problem~8.2 of Dantas--Doucha--Jung--Raunig, asking
whether the dotted implications in their summary diagram hold.  The relevant
source crop is:

\begin{center}
\includegraphics[width=\textwidth]{figures/open_problem_crop.png}
\end{center}

\section{Question, literature check, and a definition caveat}

All Banach spaces below are real.  We use the terminology of
Dantas--Doucha--Jung--Raunig \cite{DDJR}.  Their Problem 8.2 asks, among
other dotted arrows, whether
\begin{equation*}
  G\text{-}\RNP \quad\Longrightarrow\quad \RNP
  \tag{1.1}
\end{equation*}
holds for every topological group and every continuous isometric action on a
Banach space.

The arXiv record of \cite{DDJR} still contained only version 1 when this note
was prepared.  Targeted searches through June 22, 2026 found no article or
preprint resolving (1.1).  A later preprint citing \cite{DDJR} uses its
Theorem 6.4 but does not address Problem 8.2 \cite{Daugavet}.  The general
fact that exotic groups have no nontrivial nonsingular near actions is known
and is also noted explicitly in a recent paper of Avraham-Re'em and Peterzil
\cite{AP}; the connection with (1.1) below does not appear in the literature
located in this search.

There is a minor technical issue in the definitions as printed in
\cite{DDJR}.  Equation (3.2) asserts that an $X$-valued vector measure $m$
with merely $|m|\ll\mu$ induces a bounded operator
$T_m:L_1(\mu)\to X$ by $T_m\chi_A=m(A)$.  This is false without the stronger
bound $|m|\leq C\mu$.  For instance, the scalar measure with density
$t^{-1/2}$ on $[0,1]$ is absolutely continuous but does not define a bounded
functional on $L_1[0,1]$.  The natural repair is either to impose
$|m|\leq C\mu$ in the operator formulation or to change to a dominating
control measure.  The finite-group proof and the counterexample below are
valid under either repair; in the counterexample every quasi-invariant
vector measure is actually zero.

We shall use the following direction of the operator formulation, whose proof
is unaffected by this caveat.

\begin{lemma}\label{lem:operator-direction}
If a $G$-Banach space $X$ has the $G$-RNP and
$S:L_1(\Omega,\mu)\to X$ is a bounded $G$-equivariant operator for a
continuous positive isometric action of $G$ on $L_1(\mu)$, then $S$ is
representable by a function in $L_\infty(\mu;X)$.
\end{lemma}

\begin{proof}
Set $m(A)=S\chi_A$.  For every measurable $E$ and every finite measurable
partition $(E_j)$ of $E$,
\[
  \sum_j\|m(E_j)\|
  \leq \|S\|\sum_j\mu(E_j)=\|S\|\mu(E),
\]
so $|m|\leq\|S\|\mu$.  Equivariance of $S$ is exactly the required
$G$-quasi-invariance of $m$.  The $G$-RNP gives a Bochner density $F$.
Since $|m|(A)=\int_A\|F\|\,d\mu$, the variation bound implies
$\|F\|\leq\|S\|$ almost everywhere.  Thus $F\in L_\infty(\mu;X)$ and,
first on simple functions and then by continuity,
$S\varphi=\int\varphi F\,d\mu$ for every $\varphi\in L_1(\mu)$.
\end{proof}

\section{Proof intuition}

The finite-group partial works by inducing an arbitrary vector measure to
the disjoint union $G\times\Omega$ and then restricting the equivariant
density to a positive-measure sheet.  For locally compact
$\sigma$-compact groups, the same idea survives by averaging over a
quasi-invariant Haar-equivalent probability measure and then untwisting the
density.

The counterexample uses the opposite phenomenon.  Some topological groups,
including $G=\Homeo_+([0,1])$, have no nontrivial continuous Hilbert
representations.  Via the Mazur map, any admissible continuous positive
isometric $L_1$-action for such a group becomes a Hilbert representation on
$L_2$, hence must be trivial.  Therefore the $G$-RNP sees only the fixed-point
subspace $X^G$.  Taking $X=C_0((0,1))$ with the natural composition action
gives $X^G=\{0\}$, so the equivariant property holds, while $X$ contains
$c_0$ and therefore fails the classical RNP.

\section{Verification of the supplied finite-group proof}

The supplied argument is correct, including its variation estimate,
equivariance computation, and restriction to a positive-measure sheet.  For
completeness, here is a concise operator version.

\begin{proposition}\label{prop:finite}
Let $G$ be a finite Hausdorff topological group and let $X$ be a
$G$-Banach space.  If $X$ has the $G$-RNP, then $X$ has the classical RNP.
\end{proposition}

\begin{proof}
A finite Hausdorff group is discrete.  Let
$T:L_1(\Omega,\mu)\to X$ be bounded.  Give $G$ normalized counting measure
and, for $\Phi\in L_1(G\times\Omega)$, write
$\Phi_g(\omega)=\Phi(g,\omega)$.  Define
\[
  \widetilde T\Phi
  =\frac1{|G|}\sum_{g\in G}g\,T(\Phi_g).
\]
Let $G$ act on $G\times\Omega$ by left translation in the first coordinate.
Then $\widetilde T$ is bounded and $G$-equivariant.  By Lemma
\ref{lem:operator-direction}, there is
$F\in L_\infty(G\times\Omega;X)$ representing $\widetilde T$.
Fix $g_0\in G$ and test the representation formula on functions supported on
$\{g_0\}\times\Omega$.  For every $f\in L_1(\mu)$,
\[
  \frac1{|G|}g_0T(f)
  =\frac1{|G|}\int_\Omega f(\omega)F(g_0,\omega)\,d\mu(\omega).
\]
Consequently
\[
  T(f)=\int_\Omega f(\omega)
  \bigl(g_0^{-1}F(g_0,\omega)\bigr)\,d\mu(\omega).
\]
Thus every bounded operator from an $L_1$-space into $X$ is representable,
which is the classical operator characterization of the RNP.
\end{proof}

\begin{remark}
The vector-measure proof in the supplied packet is equally valid.  Its key
estimate is
\[
 |M|(E)\leq \frac1{|G|}\sum_{g\in G}|m|(E_g),
\]
which follows by applying the triangle inequality to an arbitrary finite
partition of $E$.  Since each singleton sheet has positive counting measure,
restricting the density to one sheet is legitimate.
\end{remark}

\section{Two positive extensions}

\subsection{Weak \texorpdfstring{$G$}{G}-dentability is ordinary dentability}

\begin{lemma}\label{lem:weak-collapse}
With Definition 3.5 of \cite{DDJR} exactly as printed, a $G$-Banach space is
weakly $G$-dentable if and only if it is dentable in the classical sense.
\end{lemma}

\begin{proof}
Classical dentability plainly implies weak $G$-dentability.  Conversely,
assume every bounded $G$-invariant subset of $X$ is dentable.  Let
$A\subseteq X$ be a nonempty bounded set and fix $\varepsilon>0$.  Its
saturation $G\cdot A$ is bounded and $G$-invariant, hence dentable.  Choose
$y\in G\cdot A$ such that
\[
 y\notin \overline{\co}\bigl((G\cdot A)\setminus
 B_\varepsilon(y)\bigr).
\]
Write $y=gx$ with $x\in A$.  Since $g$ is an isometry,
\[
 g\bigl(A\setminus B_\varepsilon(x)\bigr)
 \subseteq (G\cdot A)\setminus B_\varepsilon(y).
\]
If $x$ belonged to
$\overline{\co}(A\setminus B_\varepsilon(x))$, applying $g$ would
contradict the choice of $y$.  Hence $A$ is dentable.  Since $A$ and
$\varepsilon$ were arbitrary, $X$ is classically dentable.
\end{proof}

Read together with the implication $(1)\Rightarrow(3)$ in Theorem 6.4 and
Example 6.2 of \cite{DDJR}, this already gives the following consequence.

\begin{corollary}\label{cor:propertyF}
Assuming Theorem 6.4 of \cite{DDJR} under the natural definition repair
discussed above, if $G$ has property F, then $G$-RNP implies RNP.  In
particular, this holds for every locally compact $\sigma$-compact group,
without a separability assumption.
\end{corollary}

\subsection{A direct induction theorem}

The finite-sheet construction has a continuous analogue.  The fact that
individual Haar-measure sheets are null is not an obstruction: after obtaining
a density on $G\times\Omega$, one untwists and averages that density over
$G$.

\begin{theorem}\label{thm:lcsc}
Let $G$ be a locally compact $\sigma$-compact group and let $X$ be a
$G$-Banach space.  If $X$ has the $G$-RNP, then $X$ has the classical RNP.
\end{theorem}

\begin{proof}
Let $\lambda$ be a left Haar measure on $G$.  Since $G$ is
$\sigma$-compact, $\lambda$ is $\sigma$-finite.  Choose a Borel function
$w>0$ almost everywhere with $\int_Gw\,d\lambda=1$, and put
$d\nu=w\,d\lambda$.  Then $\nu$ is a probability measure equivalent to
$\lambda$.  Left translations act nonsingularly on $(G,\nu)$ and
continuously on $L_1(G,\nu)$ by positive isometries.  Indeed, the isometry
\[
 J:L_1(G,\nu)\longrightarrow L_1(G,\lambda),\qquad Jf=wf,
\]
conjugates this action to the usual strongly continuous left regular
action on $L_1(G,\lambda)$.

Let $T:L_1(\Omega,\mu)\to X$ be bounded.  On $G\times\Omega$ use the
product measure $\nu\times\mu$ and the action
$h\cdot(g,\omega)=(hg,\omega)$.  For
$\Phi\in L_1(G\times\Omega)$ define
\begin{equation*}
  \widetilde T\Phi
  =\int_G g\,T(\Phi_g)\,d\nu(g).
  \tag{3.1}
\end{equation*}
The Bochner integral is well defined.  One way to see the measurability is
that $G$ is a countable union of compact sets, so every norm-continuous orbit
in $X$ is separable; approximation of $\Phi$ by product-simple functions
then gives strong measurability.  Moreover,
\[
 \|\widetilde T\Phi\|
 \leq \|T\|\int_G\|\Phi_g\|_{L_1(\mu)}\,d\nu(g),
\]
so $\widetilde T$ is bounded.

For $h\in G$, let
\[
 r_h=\frac{d(h_*\nu)}{d\nu}.
\]
The induced positive isometry on the product $L_1$-space is
\[
 (U_h\Phi)(g,\omega)=r_h(g)\Phi(h^{-1}g,\omega).
\]
Since $r_h\,d\nu=d(h_*\nu)$, a change of variables gives
\begin{align*}
 \widetilde T(U_h\Phi)
 &=\int_G g\,r_h(g)T(\Phi_{h^{-1}g})\,d\nu(g)\\
 &=\int_G hk\,T(\Phi_k)\,d\nu(k)
 =h\,\widetilde T\Phi.
\end{align*}
Thus $\widetilde T$ is $G$-equivariant.  Lemma
\ref{lem:operator-direction} supplies
$F\in L_\infty(G\times\Omega;X)$ such that
\[
 \widetilde T\Phi
 =\int_{G\times\Omega}\Phi(g,\omega)F(g,\omega)
 \,d\nu(g)d\mu(\omega).
\]

For $A\in\Sigma$, set
\[
 q_A(g)=\int_A F(g,\omega)\,d\mu(\omega).
\]
Testing on $\chi_B(g)\chi_A(\omega)$, where $B\subseteq G$ is arbitrary and
measurable, yields
\[
 \int_Bq_A(g)\,d\nu(g)
 =\int_Bg\,T(\chi_A)\,d\nu(g).
\]
Uniqueness of Bochner-integrable functions therefore gives
\[
 q_A(g)=g\,T(\chi_A)
 \quad\text{for $\nu$-almost every }g.
\]
Define
\[
 f(\omega)=\int_Gg^{-1}F(g,\omega)\,d\nu(g).
\]
Bochner--Fubini applies and $f\in L_\infty(\Omega;X)$.  Hence, for every
$A\in\Sigma$,
\begin{align*}
 \int_Af\,d\mu
 &=\int_Gg^{-1}q_A(g)\,d\nu(g)\\
 &=\int_GT(\chi_A)\,d\nu(g)
 =T(\chi_A).
\end{align*}
It follows by linearity and continuity that
$T\varphi=\int_\Omega\varphi f\,d\mu$ for every
$\varphi\in L_1(\mu)$.  Thus $T$ is representable.  Since $T$ and the finite
measure space were arbitrary, $X$ has the RNP.
\end{proof}

This theorem contains the supplied finite-group result and also covers every
compact group and every countable discrete group.

\section{Exotic groups make the equivariant property too weak}

Call a topological group \emph{unitarily trivial} if every strongly
continuous orthogonal representation on a real Hilbert space is trivial.
Such groups are also called exotic.  The next observation is the key.

\begin{lemma}[Mazur-map transfer]\label{lem:mazur}
Let $G$ be unitarily trivial.  Every continuous action of $G$ on a finite
measure space by regular set isomorphisms is trivial modulo null sets.
Equivalently, every strongly continuous homomorphism
\[
 U:G\longrightarrow\Iso^+(L_1(\mu))
\]
is trivial.
\end{lemma}

\begin{proof}
For $f\in L_1(\mu)$ define the Mazur map
\[
 M(f)=\operatorname{sgn}(f)|f|^{1/2}\in L_2(\mu).
\]
By the positive Banach--Lamperti representation, for every $g\in G$ there
are a regular set isomorphism $\theta_g$ and a nonnegative Radon--Nikodym
weight $a_g$ such that
\[
 U_gf=a_g\,\widetilde\theta_gf.
\]
Consequently
\[
 V_g\xi=a_g^{1/2}\,\widetilde\theta_g\xi
 \qquad(\xi\in L_2(\mu))
\]
is a linear orthogonal operator and
\begin{equation*}
 M(U_gf)=V_gM(f).
 \tag{4.1}
\end{equation*}
Thus $(V_g)_{g\in G}$ is an orthogonal representation.

It is strongly continuous.  For real scalars $s,t$,
\[
 \bigl|\operatorname{sgn}(s)\sqrt{|s|}
       -\operatorname{sgn}(t)\sqrt{|t|}\bigr|^2
 \leq 2|s-t|,
\]
so
\[
 \|M(f)-M(k)\|_2^2\leq2\|f-k\|_1.
\]
For $\xi\in L_2$, put
$f=M^{-1}(\xi)=\operatorname{sgn}(\xi)|\xi|^2$.  Strong continuity of $U$
and (4.1) imply strong continuity of $V$.  Since $G$ is unitarily trivial,
$V_g=I$ for every $g$.  Injectivity of $M$ then gives $U_g=I$.
Applying this identity first to $1$ and then to characteristic functions
shows that $a_g=1$ and $\theta_gA=A$ modulo null sets for every measurable
$A$.  Thus the measure-algebra action is trivial.
\end{proof}

\begin{proposition}\label{prop:fixed}
Let $G$ be unitarily trivial and let $X$ be a $G$-Banach space.  Then
\[
 X\text{ has the }G\text{-RNP}
 \quad\Longleftrightarrow\quad
 X^G:=\{x\in X:gx=x\text{ for all }g\in G\}\text{ has the RNP}.
\]
\end{proposition}

\begin{proof}
Assume first that $X^G$ has the RNP.  By Lemma \ref{lem:mazur}, every
measure-space action occurring in the definition of the $G$-RNP is trivial.
A $G$-quasi-invariant vector measure $m$ therefore satisfies
$gm(A)=m(A)$ for every $g$ and $A$, so its range is contained in $X^G$.
The RNP of $X^G$ supplies the required density.

Conversely, assume $X$ has the $G$-RNP.  Equip an arbitrary finite measure
space with the trivial $G$-action.  Every $X^G$-valued absolutely continuous
vector measure is then $G$-quasi-invariant, hence has an $X$-valued Bochner
density $F$.  Applying the quotient map $X\to X/X^G$ to the integral
identity gives
\[
 \int_A(F+X^G)\,d\mu=0
 \qquad(A\in\Sigma).
\]
Uniqueness of Bochner densities implies $F(\omega)\in X^G$ almost
everywhere.  Hence $X^G$ has the RNP.
\end{proof}

\section{The counterexample}

Let
\[
 G=\Homeo_+([0,1]),
\]
the Polish group of orientation-preserving self-homeomorphisms of the closed
interval, with the compact-open topology.  Let
\[
 X=C_0((0,1))
 =\{f\in C[0,1]:f(0)=f(1)=0\}
\]
with the supremum norm, and define
\begin{equation*}
 (g\cdot f)(t)=f(g^{-1}(t)).
 \tag{5.1}
\end{equation*}

\begin{theorem}\label{thm:counterexample}
The action in (5.1) is a continuous faithful action by linear isometries,
$X$ has the $G$-RNP, and $X$ does not have the classical RNP.  Consequently,
the implication $G$-RNP $\Rightarrow$ RNP is false for general topological
groups.
\end{theorem}

\begin{proof}
Each operator in (5.1) is a linear isometry.  If $g_i\to g$ in the
compact-open topology, then $g_i^{-1}\to g^{-1}$ uniformly.  Since every
$f\in X$ is uniformly continuous on $[0,1]$,
\[
 \|g_i\cdot f-g\cdot f\|_\infty\longrightarrow0.
\]
Thus the representation is strongly continuous, and hence the action is
jointly continuous.  It is faithful because continuous functions separate
points.

Megrelishvili proved that every weakly continuous representation of
$\Homeo_+([0,1])$ by linear isometries on a reflexive Banach space is trivial
\cite{Megrelishvili}.  In particular, $G$ is unitarily trivial.

Moreover, $X^G=\{0\}$.  The group $G$ acts transitively on $(0,1)$, so an
invariant function is constant there; continuity and the endpoint conditions
force that constant to be zero.  Proposition \ref{prop:fixed} now gives the
$G$-RNP of $X$.

It remains to show that $X$ does not have the RNP.  Choose pairwise disjoint
closed intervals $I_n\subset(0,1)$ accumulating only at $0$, and choose tent
functions $\varphi_n\in X$ with
$\supp\varphi_n\subseteq I_n$ and $\|\varphi_n\|_\infty=1$.  Then
\[
 J:c_0\longrightarrow X,
 \qquad J((a_n))=\sum_{n=1}^\infty a_n\varphi_n,
\]
is a linear isometry.  Indeed, the supports are disjoint, the supremum norm
is $\sup_n|a_n|$, and continuity at the accumulation point follows from
$a_n\to0$.  Its range is closed.  Since the RNP passes to closed subspaces
and $c_0$ fails the RNP, $X$ fails the RNP.
\end{proof}

For completeness, the last standard fact has a short direct proof.

\begin{lemma}\label{lem:c0}
The Banach space $c_0$ does not have the RNP.
\end{lemma}

\begin{proof}
Let $(r_n)$ be the Rademacher functions on $[0,1]$ and define
\[
 m(A)=\left(\int_A r_n\,d\lambda\right)_{n\geq1}.
\]
For every measurable $A$, Bessel's inequality applied to
$\chi_A\in L_2[0,1]$ shows that the coordinates tend to zero, so
$m(A)\in c_0$.  Also $\|m(A)\|_\infty\leq\lambda(A)$.  This estimate shows
countable additivity in the $c_0$ norm and, for every finite measurable
partition $(A_j)$ of $A$,
\[
 \sum_j\|m(A_j)\|_\infty\leq\sum_j\lambda(A_j)=\lambda(A),
\]
so $|m|\leq\lambda$.

If $m$ had a Bochner density $F:[0,1]\to c_0$, applying the $n$th coordinate
functional would give $F_n=r_n$ almost everywhere.  Outside one common null
set this holds for every $n$, so
$F(t)=(r_n(t))_{n\geq1}$.  Every coordinate has absolute value $1$, and this
sequence is not in $c_0$, a contradiction.
\end{proof}

\section{Conclusion}

The supplied finite-group result is valid.  The positive statement extends
to all locally compact $\sigma$-compact groups, both by the direct induction
argument above and, after observing Lemma \ref{lem:weak-collapse}, through
the source paper's property-F theorem.  The unrestricted implication fails:
for exotic groups the admissible measure-algebra actions can all be trivial,
so the $G$-RNP sees only the fixed-point subspace $X^G$ and can miss the
geometry of the rest of $X$ completely.

\section{Scope and human review recommendation}

The packet is a candidate counterexample to the unrestricted dotted
implication $G$-RNP $\Rightarrow$ RNP in Problem~8.2.  It is not a negative
answer for locally compact $\sigma$-compact groups; the packet proves the
positive implication in that class.  Human review should focus on the
definition repair in Section~2, the Mazur-map transfer from positive
$L_1$-actions to Hilbert representations, the exact form of Megrelishvili's
triviality theorem, and whether the source paper's intended scope excludes
non-locally-compact topological groups.

\begin{thebibliography}{9}

\bibitem{DDJR}
S. Dantas, M. Doucha, M. Jung and T. Raunig,
\emph{Group equivariant Radon--Nikodym property and its characterizations},
arXiv:2510.23100v1, 2025.

\bibitem{Daugavet}
S. Dantas, H. del R\'io and T. Raunig,
\emph{A group action approach to the Daugavet property},
arXiv:2606.20429v1, 2026.

\bibitem{AP}
N. Avraham-Re'em and G. Peterzil,
\emph{Wasserstein stability and the nonsingular Borel lifting problem},
arXiv:2606.03721v1, 2026.

\bibitem{Megrelishvili}
M. G. Megrelishvili,
\emph{Every semitopological semigroup compactification of the group
$H_+[0,1]$ is trivial},
Semigroup Forum \textbf{63} (2001), 357--370.
\href{https://doi.org/10.1007/s002330010076}{doi:10.1007/s002330010076}.

\end{thebibliography}

\end{document}
